\documentclass[a4paper,12pt]{article}
\usepackage[utf8]{inputenc}
\usepackage[spanish]{babel}
\usepackage{amsmath, amssymb, amsthm}
\usepackage{geometry}
\usepackage{booktabs}
\usepackage{enumitem}
\usepackage{tcolorbox}
\usepackage{multicol}

% Configuración de márgenes
\geometry{top=2.5cm, bottom=2.5cm, left=2.5cm, right=2.5cm}

% Personalización de títulos
\title{\textbf{Resumen Técnico: Series Infinitas}\\ \large Basado en J. Stewart - Capítulo 11}
\author{Cálculo de Trascendentes Tempranas}
\date{}

\begin{document}

\maketitle
\section{Series Infinitas}

\subsection{Definiciones Fundamentales}

\subsubsection{Definición de Serie}
Una serie infinita es la suma de los términos de una sucesión infinita $\{a_n\}_{n=1}^{\infty}$ y se denota como:
$$\sum_{n=1}^{\infty} a_n = a_1 + a_2 + a_3 + \dots + a_n + \dots$$

\subsubsection{Sumas Parciales y Convergencia}
Para determinar si una serie tiene una suma finita, se analiza la \textbf{sucesión de sumas parciales} $\{s_n\}$:
\begin{itemize}
    \item $s_1 = a_1$
    \item $s_2 = a_1 + a_2$
    \item $s_n = \sum_{i=1}^{n} a_i$
\end{itemize}

\begin{tcolorbox}[title=Definición de Convergencia,colback=blue!5!white,colframe=blue!75!black]
Si la sucesión $\{s_n\}$ es convergente y $\lim_{n \to \infty} s_n = s$ (donde $s$ es un número real), la serie $\sum a_n$ es \textbf{convergente} y su suma es $s$. Si el límite no existe o es infinito, la serie es \textbf{divergente}.
\end{tcolorbox}

\subsection{Propiedades}
\begin{itemize}
    \item La convergencia o divergencia de una serie no es afectada:
    \begin{itemize}
        \item al cambiar un número finito de términos.
        \item al multiplicar todos los términos por una constante: $\sum_{n=1}^{\infty} a_n  C = C \sum_{n=1}^{\infty} a_n$
    \end{itemize} 
    \item Dos series convergentes pueden sumar o restarse término a término y la serie resultante también converge y su suma $s$ es obtenida sumando o restando las sumas de las series dadas:
    $$\sum_{n=1}^{\infty} a_n \pm \sum_{n=1}^{\infty} b_n = \sum_{n=1}^{\infty} (a_n \pm b_n)$$
    \item Si se suma una serie convergente y una divergente, la serie resultante es divergente.
\end{itemize}

\subsection{Series Especiales}

\subsubsection{Serie Geométrica}
Se define como $\sum_{n=1}^{\infty} ar^{n-1}$. 
\begin{itemize}
    \item \textbf{Converge} si $|r| < 1$, su suma es $s = \frac{a}{1-r}$.
    \item \textbf{Diverge} si $|r| \geqslant 1$.
\end{itemize}
\textit{Resumen:} la suma de una serie geométrica convergente es: $\frac{\text{primer término}}{1 - \text{razón común}}$

\subsubsection{Serie p}
Se define como $\sum_{n=1}^{\infty} \frac{1}{n^p}$.
\begin{itemize}
    \item \textbf{Converge} si $p > 1$.
    \item \textbf{Diverge} si $p \leqslant 1$.
\end{itemize}

\subsubsection{Serie Armónica}
Es el caso especial de la serie $p$ donde $p=1$: $\sum_{n=1}^{\infty} \frac{1}{n}$.
\textbf{Siempre es divergente.}

\subsection{Pruebas de Convergencia}

\subsubsection{Prueba de la Divergencia (Término n-ésimo)}
Si $\lim_{n \to \infty} a_n$ no existe o si $\lim_{n \to \infty} a_n \neq 0$, entonces la serie $\sum a_n$ es \textbf{divergente}.
\begin{tcolorbox}
    Si $\sum_{n=1}^{\infty} a_n$ es \textbf{convergente}, entonces $\lim_{n \to \infty} a_n = 0$, \textit{\textbf{pero el recíproco no es necesariamente cierto}}.
\end{tcolorbox}







\subsubsection{Prueba de la Integral}
Se aplica si $a_n = f(n)$, donde $f$ es continua, positiva y decreciente en $[1, \infty)$. La serie $\sum a_n$ converge si y solo si la integral impropia $\int_{1}^{\infty} f(x) dx$ converge.
\textit{Ejemplo:} $\sum \frac{1}{n^2+1}$. La integral $\int_{1}^{\infty} \frac{1}{x^2+1} dx = \pi/4$, por lo que la serie converge.

\subsubsection{Pruebas de Comparación}
\begin{enumerate}
    \item \textbf{Comparación Directa:} Si $0 \leq a_n \leq b_n$ y $\sum b_n$ converge $\implies \sum a_n$ converge. Si $a_n \geq b_n \geq 0$ y $\sum b_n$ diverge $\implies \sum a_n$ diverge.
    \item \textbf{Comparación del Límite:} Si $\lim_{n \to \infty} \frac{a_n}{b_n} = c$ (donde $c > 0$ y es finito), ambas series tienen el mismo comportamiento.
\end{enumerate}

\subsubsection{Prueba de la Serie Alternante}
Para series de la forma $\sum (-1)^{n-1} b_n$ con $b_n > 0$. Converge si:
\begin{itemize}
    \item $b_{n+1} \leq b_n$ para toda $n$.
    \item $\lim_{n \to \infty} b_n = 0$.
\end{itemize}

\subsubsection{Prueba de la Razón (D'Alembert)}
Sea $L = \lim_{n \to \infty} | \frac{a_{n+1}}{a_n} |$:
\begin{itemize}
    \item Si $L < 1$, es \textbf{absolutamente convergente}.
    \item Si $L > 1$, es \textbf{divergente}.
    \item Si $L = 1$, la prueba no concluye.
\end{itemize}

\subsubsection{Prueba de la Raíz (Cauchy)}
Sea $L = \lim_{n \to \infty} \sqrt[n]{|a_n|}$:
\begin{itemize}
    \item Si $L < 1$, converge absolutamente; si $L > 1$, diverge.
\end{itemize}

\subsection{Convergencia Absoluta y Condicional}
\begin{itemize}
    \item \textbf{Convergencia Absoluta:} $\sum a_n$ es absolutamente convergente si $\sum |a_n|$ es convergente.
    \item \textbf{Convergencia Condicional:} Una serie es condicionalmente convergente si es convergente pero $\sum |a_n|$ diverge.
\end{itemize}

\subsection{Series de Potencias}
Es una serie de la forma $\sum_{n=0}^{\infty} c_n (x - a)^n$.
\begin{itemize}
    \item \textbf{Radio de convergencia ($R$):} Distancia desde el centro $a$ donde la serie converge.
    \item \textbf{Intervalo de convergencia:} Conjunto de valores de $x$ para los cuales la serie converge.
\end{itemize}

\subsection{Series de Taylor y Maclaurin}
Si $f$ tiene representación en serie de potencias centrada en $a$:
$$ f(x) = \sum_{n=0}^{\infty} \frac{f^{(n)}(a)}{n!} (x-a)^n $$
\textbf{Serie de Maclaurin:} Caso donde $a = 0$.

\subsubsection{Resto de Taylor}
Se define como $R_n(x) = f(x) - T_n(x)$. Si $\lim_{n \to \infty} R_n(x) = 0$, entonces la serie converge a la función $f(x)$.

\end{document}